% Noether P09 Korean producer draft — UNCHECKED
% T03-U20; ED0002 lines6573--6594; source 1,121 LF B / 795975CAC70B5A2C43F8C8A94D03C5B3DF7FF1C8095A706AE3BA663361B7862F
% Pointer v009 / ED0002: B06BE3530D9CF2E82B56FDBA7FE41D5D044DF2425DFA2A059D4939EAA2F7A6C2 / C9A125167ACB33D914EE4374B65AE7CDF0052F568371B8B77B720EA178ABF0E3
% Translation only; no review, build, render, or approval.

\emph{그러나 어떠한 정렬에서도 $\mathfrak G$로부터, $\mathfrak G$가 기저의 수들의 모든 분수형 유리함수와 분수형 대수함수를 배제하도록 기저를 골라낼 수는 없다.} 실제로 영역의 어떤 함수 $z(\eta)$를 기저의 수로 고른다고 하자. 이 함수는 다음 방정식으로 정의된다.
\[
 z^k+A_1(\eta)z^{k-1}+\cdots+A_k(\eta)=0,
\]
여기서
\[
 A_i(\eta)=a_i\cdot\frac{E_1(\eta)}{E_2(\eta)}
\]
는 정수적 유리 함수식이다. 그러면 함수 $\frac{a_k}{z}$는 $\mathfrak G$에 속하며, $\frac{a_k}{\sqrt z}$, $\frac{a_k}{\sqrt[3]z}$ 등도 마찬가지이다.\srcfn{*)}{완전히 같은 방식으로, $\eta$의 모든 대수함수에 어떤 유리정수를 곱하여 함수식 영역 $\mathfrak G$의 원소로 만들 수 있다. 따라서 이 영역은 \emph{모든 복소수를 $\mathfrak G$의 한 원소와 한 유리정수의 몫으로 나타낼 수 있다}는 성질을 갖는다. 이는 대수적 수의 경우와 유사하다.} 따라서 \emph{영역 $\mathfrak G$ 전체에 대하여 기저 성질 4)와 5)가 배제된다.}
